% macros.tex -- shared notation for all sections
\newtheorem{theorem}{Theorem}[section]
\newtheorem{lemma}[theorem]{Lemma}
\newtheorem{proposition}[theorem]{Proposition}
\newtheorem{corollary}[theorem]{Corollary}
\newtheorem{fact}[theorem]{Computed Fact}
\theoremstyle{definition}
\newtheorem{definition}[theorem]{Definition}
\newtheorem{example}[theorem]{Example}
\theoremstyle{remark}
\newtheorem{remark}[theorem]{Remark}
% players and stones
\newcommand{\Maker}{Maker}            % first player
\newcommand{\Breaker}{Breaker}        % second player
\newcommand{\Q}{Q}                    % the board {0,...,5}^3 (216 cells)
\newcommand{\Lines}{\mathcal{L}}      % the 148 lines
\newcommand{\live}[1]{\mathcal{L}(#1)}% live lines w.r.t. Breaker's stones #1
\newcommand{\N}{N}                    % N(S) = union of the lines in S
\newcommand{\sur}{\operatorname{sur}} % surplus  sur_F(S) = |N(S)\F| - 2|S|
\newcommand{\dfc}{\operatorname{def}} % deficiency def_Y(S) = 2|S| - |N(S)\Y|
\newcommand{\slack}{\operatorname{slack}}
\newcommand{\G}{G}                    % symmetry group, |G| = 1152
\newcommand{\canon}{\operatorname{canon}}
\newcommand{\antip}[1]{\overline{#1}} % antipode
\newcommand{\Busy}{\mathcal{K}}       % the 24 busy cells
\newcommand{\CC}{\mathsf{C}}          % C(B,F,k): pairing condition
\newcommand{\CPhi}{\mathsf{C}_\Phi}   % C_Phi(B,F,k): surplus condition
\newcommand{\LEAF}{\mathsf{LEAF}}     % LEAF(B,F,k)
\newcommand{\WON}{\textup{\textsc{won}}}
\newcommand{\VOID}{\textup{\textsc{void}}}
\newcommand{\LOST}{\mathsf{LOST}}
\newcommand{\Pst}{\mathsf{P}}         % P-state (Maker to move): \Pst(B,F,t)
\newcommand{\Qst}{\mathsf{Q}}         % Q-state (Breaker to move): \Qst(B,F,t)
\newcommand{\Rmax}{R_{\max}}
\newcommand{\Tbl}{\mathrm{Tbl}}
\newcommand{\Lv}{\mathrm{Lv}}
\newcommand{\Tseven}{T_7}
\newcommand{\Teight}{T_8}
\newcommand{\computed}{\textup{\textsc{computed}}}
\newcommand{\code}[1]{\texttt{#1}}
% --- Section 5 (abstract game) ---
\newcommand{\Etwo}{\mathrm{E2}}        % pattern E2 (a threat and a danger pair)
\newcommand{\Ethree}{\mathrm{E3}}      % pattern E3 (two danger pairs)
\newcommand{\TABLE}{\textup{\textsc{table}}}    % TABLE mode of sigma
\newcommand{\PAIRING}{\textup{\textsc{pairing}}}% PAIRING mode of sigma
\newcommand{\WONp}{\WON^{+}}           % WON in the eager variant
\newcommand{\rk}{\operatorname{rk}}    % rank of a state (not \rho, which is the reversal i -> 5-i)
% --- computational sections ---
\DeclareRobustCommand{\fn}[1]{{\upshape\path{#1}}} % file names and command-line options (robust; always upright)
% provenance paragraph of a Computed Fact: the statement is in italics, the provenance upright
\newcommand{\established}{\upshape\emph{Established by:}\space}
\newcommand{\accept}{\textup{\textsc{accept}}}
\newcommand{\reject}{\textup{\textsc{reject}}}
\newcommand{\worst}{\operatorname{worst}} % worst(B,F,k): largest violation of Hall's condition
% --- Section 9 (the 7x7x7 board) ---
\newcommand{\Qseven}{\Q_7}             % the board {0,...,6}^3
% --- cross-reference names (cleveref) ---
\crefname{fact}{Computed Fact}{Computed Facts}
\Crefname{fact}{Computed Fact}{Computed Facts}
\crefname{algocf}{Algorithm}{Algorithms}
\Crefname{algocf}{Algorithm}{Algorithms}
\crefname{appendix}{Appendix}{Appendices}
\Crefname{appendix}{Appendix}{Appendices}

% --- revision: notation clashes reported by the journal referee (N4) ---
% The board {0,...,5}^3 is written \Omega: the letter Q is reserved for the Q-states \Qst(B,F,t) and
% for the sets Q_a, Q_b of revealed stones in the danger children.  \Qseven = \Q_7 follows.
\renewcommand{\Q}{\Omega}

% --- revision, second phase ---
% \needlines{n}: start a new page here unless about n lines are left on the current page
% (used before the main theorem, which must not be broken across pages)
\newcommand{\needlines}[1]{\par\penalty-100\begingroup\dimen0=#1\baselineskip
  \vskip 0pt plus \dimen0\penalty-200\vskip 0pt plus -\dimen0\endgroup}
% size of the supplementary archive as stated in Section 8.1; TO BE UPDATED whenever the archive is
% rebuilt (scripts/make_artifact.sh prints the numbers and stops if \ArchiveFiles is not the number of
% files in the archive, MANIFEST.sha256 included)
\newcommand{\ArchiveSize}{about $6.3$\,MB}
\newcommand{\ArchiveUnpacked}{about $38$\,MB}
\newcommand{\ArchiveFiles}{355}
% captions of longtables (Tables of Sections 7.4 and 8.1) in the style of the amsart table captions
\makeatletter
\@ifundefined{LT@makecaption}{}{%
\def\LT@makecaption#1#2#3{%
  \LT@mcol\LT@cols c{\hbox to\z@{\hss\parbox[t]{\dimexpr\textwidth-2\captionindent\relax}{%
    \normalfont\normalsize
    \setbox\@tempboxa\hbox{{\scshape #2.}\enspace #3}%
    \ifdim\wd\@tempboxa>\hsize \noindent{\scshape #2.}\enspace #3\par
    \else \centering{\scshape #2.}\enspace #3\par\fi}\hss}}}}
\makeatother
% file names (\fn = \path) are broken only after a slash or a comma, not inside a name
% (journal referee: "spotcheck_ full.json", "ip_regression. json")
\renewcommand{\UrlBreaks}{\do\/\do\,}
